% Part: first-order-logic
% Chapter: tableaux
% Section: proof-theoretic-notions

\documentclass[../../../include/open-logic-section]{subfiles}

\begin{document}

\iftag{FOL}
      {\olfileid{fol}{tab}{ptn}}
      {\olfileid{pl}{tab}{ptn}}

\olsection{నిరూపణ-సిద్ధాంత భావనలు}

\begin{editorial}
ఈ విభాగం టాబ్లోలకు నిరూప్యతా సంబంధం, అవైరుధ్యం యొక్క నిర్వచనాలను
సమీకరిస్తుంది.
\end{editorial}

\begin{explain}
అనేక ముఖ్యమైన అర్థపర భావనలను (చెల్లుబాటు, అనుగమనం, సంతృప్తత)
నిర్వచించినట్లే, ఇప్పుడు వాటికి అనుగుణమైన \emph{నిరూపణ-సిద్ధాంత
భావనలను} నిర్వచిస్తాం. వీటిని \tetoken{నిర్మాణాలలో}{!!{structure}s} \tetoken{వాక్యాలు}{!!{sentence}s}
సంతృప్తి చెందడాన్ని ఆశ్రయించి నిర్వచించం; కొన్ని సంవృత
\tetoken{టాబ్లో}{!!{tableau}}x ఉనికిని ఆశ్రయించి నిర్వచిస్తాం. ఈ భావనలు ఒకటే అవుతాయని
కనుగొనడం ముఖ్యమైన విజయం. అవి ఒకటే అవుతాయనేదే \emph{నిర్దుష్టత},
\emph{సంపూర్ణత} సిద్ధాంతాల విషయం.
\end{explain}


\begin{defn}[సిద్ధాంతాలు]
$\sFmla{\False}{!A}$ కోసం ఒక సంవృత \tetoken{టాబ్లో}{!!{tableau}} ఉంటే
\tetoken{వాక్యం}{!!{sentence}}~$!A$ను \emph{సిద్ధాంతం} అంటాం. $!A$ ఒక సిద్ధాంతమైతే
$\Proves !A$ అని, కాకపోతే $\Proves/ !A$ అని రాస్తాం.
\end{defn}

\begin{defn}[\tetoken{వ్యుత్పాద్యత}{!!^{derivability}}]
\tetoken{వాక్యాలు}{!!{sentence}s} సమితి~$\Gamma$ నుంచి \tetoken{ఒక వాక్యం}{!!^a{sentence}}~$!A$
\emph{\tetoken{వ్యుత్పాదించదగిన}{!!{derivable}}}, అంటే $\Gamma \Proves !A$, అవుతుంది అంటే,
$\{!B_1, \dots, !B_n\} \subseteq \Gamma$ అనే పరిమిత సమితి ఒకటి,
కింది సమితికి ఒక సంవృత \tetoken{టాబ్లో}{!!{tableau}} ఉండాలి:
\[
\{
  \sFmla{\False}{!A},
  \sFmla{\True}{!B_1}, \dots,
  \sFmla{\True}{!B_n}
  \}.
\]
$\Gamma$ నుంచి $!A$ \tetoken{వ్యుత్పాదించదగిన}{!!{derivable}} కాకపోతే $\Gamma \Proves/ !A$
అని రాస్తాం.
\end{defn}

\begin{defn}[అవైరుధ్యం]
\tetoken{వాక్యాలు}{!!{sentence}s} సమితి~$\Gamma$ \emph{వైరుధ్యభరితమైనది} అంటే,
$\{!B_1, \dots, !B_n\} \subseteq \Gamma$ అనే పరిమిత సమితి ఒకటి,
కింది సమితికి ఒక సంవృత \tetoken{టాబ్లో}{!!{tableau}} ఉండాలి:
\[
\{
  \sFmla{\True}{!B_1}, \dots,
  \sFmla{\True}{!B_n}
  \}.
\]
$\Gamma$ వైరుధ్యభరితం కాకపోతే దాన్ని \emph{అవైరుధ్యమైనది} అంటాం.
\end{defn}

\begin{prop}[స్వావర్తనత్వం]
\ollabel{prop:reflexivity}
$!A \in \Gamma$ అయితే $\Gamma \Proves !A$.
\end{prop}

\begin{proof}
$!A \in \Gamma$ అయితే $\{!A\}$ అనేది $\Gamma$ యొక్క పరిమిత
ఉపసమితి; కింది \tetoken{టాబ్లో}{!!{tableau}} సంవృతమైనది:
\begin{oltableau}
  [\sFmla{\False}{\formula{A}}, just = \TAss
    [\sFmla{\True}{\formula{A}}, just = \TAss,close]
  ]
\end{oltableau}
\end{proof}


\begin{prop}[ఏకదిశత]
\ollabel{prop:monotonicity}
$\Gamma \subseteq \Delta$, $\Gamma \Proves !A$ అయితే $\Delta
\Proves !A$.
\end{prop}

\begin{proof}
$\Gamma$ యొక్క ప్రతి పరిమిత ఉపసమితీ $\Delta$కు కూడా పరిమిత ఉపసమితి.
\end{proof}

\begin{prop}[సంక్రామకత్వం]
\ollabel{prop:transitivity}
$\Gamma \Proves !A$, $\{!A\} \cup
\Delta \Proves !B$ అయితే $\Gamma \cup \Delta \Proves !B$.
\end{prop}

\begin{proof}
$\{!A\} \cup \Delta \Proves !B$ అయితే, $\Delta_0 =
\{!C_1, \dots, !C_n\} \subseteq \Delta$ అనే పరిమిత ఉపసమితి ఉండి,
\begin{align*}
\{\sFmla{\False}{!B}, & \sFmla{\True}{!A}, \sFmla{\True}{!C_1},
\dots, \sFmla{\True}{!C_n}\}
\intertext{కు ఒక సంవృత \tetoken{టాబ్లో}{!!{tableau}} ఉంటుంది. $\Gamma \Proves !A$
  అయితే, $\{!D_1, \dots, !D_m\} \subseteq \Gamma$ అయ్యేలా}
\{\sFmla{\False}{!A}, & \sFmla{\True}{!D_1},
\dots, \sFmla{\True}{!D_m}\}
\end{align*}
కు ఒక సంవృత \tetoken{టాబ్లో}{!!{tableau}} ఉంటుంది.
\sourcecorrection{OLTETAB-004}{ఈ చోట ఆంగ్ల మూలం వ్యక్తిగత
$!D_1$, \dots, $!D_m$లను నేరుగా $\Gamma$కు “ఉపసమితి”గా రాసింది.
నిగమ్యత నిర్వచనం, వెంటనే ఉపయోగించే పరిమిత పరికల్పనల సమితి ప్రకారం
సమితి కుండలీకరణాలను చేర్చి
$\{!D_1, \dots, !D_m\} \subseteq \Gamma$గా సరిచేశాం.}

ఇప్పుడు కింది పరికల్పనలు గల \tetoken{టాబ్లోను}{!!{tableau}} పరిశీలించండి:
\[
\sFmla{\False}{!B},
\sFmla{\True}{!C_1}, \dots, \sFmla{\True}{!C_n},
\sFmla{\True}{!D_1}, \dots, \sFmla{\True}{!D_m}.
\]
$!A$పై \Cut{} నియమాన్ని వర్తింపజేయండి. దాని వల్ల రెండు శాఖలు
వస్తాయి: ఒకదానిలో $\sFmla{\True}{!A}$, మరొకదానిలో
$\sFmla{\False}{!A}$ ఉంటాయి. అప్పుడు ఒక శాఖలో కిందివన్నీ లభిస్తాయి:
\[
\{\sFmla{\False}{!B}, \sFmla{\True}{!A},
\sFmla{\True}{!C_1}, \dots, \sFmla{\True}{!C_n}\}
\]
ఈ పరికల్పనలకు సంవృత \tetoken{టాబ్లో}{!!{tableau}} ఉంది కనుక దాన్ని ఆ శాఖకు
జోడించవచ్చు; $\sFmla{\True}{!A}$ గుండా వెళ్లే ప్రతి శాఖా
సంవృతమవుతుంది. మరొక శాఖలో కిందివన్నీ లభిస్తాయి:
\[
\{\sFmla{\False}{!A}, \sFmla{\True}{!D_1}, \dots,
\sFmla{\True}{!D_m}\}
\]
కనుక మరొక వైపునూ పూర్తిచేసి ఒక సంవృత \tetoken{టాబ్లోను}{!!{tableau}} పొందవచ్చు. ఇది
$\Gamma \cup \Delta \Proves !B$ అని చూపుతుంది.
\end{proof}

ముఖ్యంగా, $\Gamma \Proves !A$, $!A \Proves !B$ అయితే $\Gamma
\Proves !B$ అని దీని అర్థమని గమనించండి. అలాగే $!A_1,
\dots, !A_n \Proves !B$ అయి, ప్రతి~$i$కి $\Gamma \Proves !A_i$
అయితే $\Gamma \Proves !B$ అని కూడా అనుగమిస్తుంది.

\begin{prop}
\ollabel{prop:incons}
$\Gamma$ వైరుధ్యభరితమైనది అయితే, అప్పుడు మాత్రమే ప్రతి
\tetoken{వాక్యం}{!!{sentence}}~$!A$కు $\Gamma \Proves !A$.
\end{prop}

\begin{proof}
అభ్యాసం.
\end{proof}

\tagprob{FOL}
\begin{prob}
\olref[fol][tab][ptn]{prop:incons}ను నిరూపించండి.
\end{prob}
\tagendprob

\tagprob{notFOL}
\begin{prob}
\olref[pl][tab][ptn]{prop:incons}ను నిరూపించండి.
\end{prob}
\tagendprob

\begin{prop}[సంహతత్వం]
  \ollabel{prop:proves-compact}
  \begin{enumerate}
  \item $\Gamma \Proves !A$ అయితే, $\Gamma_0 \Proves !A$ అయ్యే
    పరిమిత ఉపసమితి $\Gamma_0 \subseteq \Gamma$ ఒకటి ఉంటుంది.
  \item $\Gamma$లోని ప్రతి పరిమిత ఉపసమితీ అవైరుధ్యమైనదైతే,
    $\Gamma$ అవైరుధ్యమైనది.
  \end{enumerate}
\end{prop}

\begin{proof}
  \begin{enumerate}
    \item $\Gamma \Proves !A$ అయితే, పరిమిత ఉపసమితి
      $\Gamma_0 = \{!B_1, \dots, !B_n\}$కు కింది సమితిపై ఒక సంవృత
      \tetoken{టాబ్లో}{!!{tableau}} ఉంటుంది:
      \[
      \{\sFmla{\False}{!A}, \sFmla{\True}{!B_1}, \dots, \sFmla{\True}{!B_n}\}
      \]
      ఇదే \tetoken{టాబ్లో}{!!{tableau}}, $\Gamma_0 \Proves !A$ అని కూడా చూపుతుంది.
    \item $\Gamma$ వైరుధ్యభరితమైతే, ఏదో పరిమిత ఉపసమితి
      $\Gamma_0 = \{!B_1, \dots, !B_n\}$ కోసం కింది సమితికి ఒక
      సంవృత \tetoken{టాబ్లో}{!!{tableau}} ఉంటుంది:
      \[
      \{\sFmla{\True}{!B_1}, \dots, \sFmla{\True}{!B_n}\}
      \]
      ఈ సంవృత \tetoken{టాబ్లో}{!!{tableau}}, $\Gamma_0$ వైరుధ్యభరితమని చూపుతుంది.
  \end{enumerate}
\end{proof}

\end{document}
